\documentclass[10pt,a4paper]{article}
\usepackage[utf8]{inputenc}
\usepackage[T1]{fontenc}
\usepackage[scaled=0.92]{DejaVuSerif}
\usepackage[scaled=0.90]{DejaVuSans}
\usepackage[scaled=0.88]{DejaVuSansMono}
\usepackage[margin=1.6cm,top=1.3cm,bottom=1.3cm]{geometry}
\usepackage{enumitem}
\usepackage{booktabs}
\usepackage{titlesec}
\usepackage{xcolor}

\definecolor{darkblue}{RGB}{0, 68, 136}

\titleformat*{\subsection}{\normalsize\bfseries}
\titlespacing*{\subsection}{0pt}{9pt}{2pt}

\pagestyle{empty} % Clean 1-page quick reference

\begin{document}

{\Large\bfseries\color{darkblue} 3D Slicer Bounding Box Navigator} \hfill {\small\bfseries Quick User Guide}\\[-0.2cm]
\rule{\textwidth}{0.8pt}

\vspace{0.05cm}
\subsection*{First-Time Setup}
\begin{enumerate}[leftmargin=1.5em, itemsep=1pt, topsep=1.5pt]
    \item \textbf{Extract the ZIP}: Right-click the downloaded zip file $\rightarrow$ \textbf{Extract All...}
    \item Open \textbf{3D Slicer}.
    \item In the menu, go to \textbf{Edit} $\rightarrow$ \textbf{Application Settings} $\rightarrow$ \textbf{Modules}.
    \item Under \textbf{Additional module paths}, click \textbf{Add} and select the \textbf{\texttt{BoundingBoxNavigator}} folder.
    \item \textit{(Recommended)} Set \textbf{Default module} to \textbf{\texttt{BoundingBoxNavigator}} so it opens automatically on startup.
    \item Click \textbf{OK} and restart Slicer.
\end{enumerate}

\subsection*{How to Annotate}
\begin{enumerate}[leftmargin=1.5em, itemsep=3pt, topsep=1.5pt]
    \item \textbf{Open the Module (if panel is not visible)}
    \begin{itemize}[leftmargin=1.2em, itemsep=0.5pt]
        \item If not opened automatically on startup, go to: \textbf{Modules} $\rightarrow$ \textbf{Annotation} $\rightarrow$ \textbf{BoundingBoxNavigator} in the top toolbar to display the panel on the right.
    \end{itemize}

    \item \textbf{Select Folders \& Scan}
    \begin{itemize}[leftmargin=1.2em, itemsep=0.5pt]
        \item Set \textbf{Input Scans Folder} (where your CT scans are) and \textbf{Output Annotations Folder}.
        \item Click \textbf{Scan Folder}. The first unreviewed case loads automatically.
        \item Each slice view shows the \textbf{original} scan for that plane when the file exists; otherwise it reconstructs from the axial volume (a warning is shown).
    \end{itemize}

    \item \textbf{Draw Bounding Boxes}
    \begin{itemize}[leftmargin=1.2em, itemsep=0.5pt]
        \item Choose a \textbf{Finding type} (default: PM nodule). Then press \textbf{\fbox{\textbf{B}}} (or click \textbf{+ Add Bounding Box}).
        \item \textbf{Click and drag} across the lesion on any slice view (axial, coronal, or sagittal).
        \item Drag the colored handles on the box to resize it in width, height, and depth.
        \item Scroll through slices with your mouse wheel to check coverage.
        \item Use the large \textbf{slice thickness} buttons (3\,mm, TS, \ldots). For TS (or any thickness without native COR/SAG), coronal and sagittal show reconstructions of that axial scan.
        \item Press \textbf{\fbox{\textbf{B}}} again for each additional lesion.
        \item \textit{To remove a box}: click it in the list and press \textbf{Delete Selected}.
    \end{itemize}

    \item \textbf{Notes / PCI Score}
    \begin{itemize}[leftmargin=1.2em, itemsep=0.5pt]
        \item Type the PCI score or clinical comments in the \textbf{Notes / PCI} field (optional).
        \item Optional: in \textbf{4.\ Segmentation}, load a \texttt{.seg.nrrd} or \texttt{.nii.gz} overlay (\texttt{.nii.gz} is imported as a labelmap, not a CT volume).
    \end{itemize}

    \item \textbf{Save and Next Scan}
    \begin{itemize}[leftmargin=1.2em, itemsep=0.5pt]
        \item Click \textbf{Save + Next} (large green button).
        \item Each box is saved as \texttt{R\_1.mrk.json}, \texttt{A\_1.mrk.json}, \ldots inside a folder named after the case ID.
        \item The next unreviewed scan loads immediately.
        \item \textit{If there are 0 lesions (PCI 0)}: click \textbf{Save + Next} and confirm with \textbf{Yes}.
    \end{itemize}
\end{enumerate}

\subsection*{Mouse \& Keyboard Shortcuts}
\begin{center}
\small
\renewcommand{\arraystretch}{1.18}
\begin{tabular}{ll}
\toprule
\textbf{Shortcut / Control} & \textbf{Action} \\
\midrule
\textbf{B} & Draw a new bounding box \\
\textbf{Mouse Wheel} & Scroll through CT slices \\
\textbf{Middle Click + Drag} & Pan view \\
\textbf{Right Click + Drag} & Zoom in / out \\
\textbf{Click row in box table} & Centers all slice views on that lesion \\
\textbf{Skip completed (checkbox)} & Checked by default. Uncheck to browse or edit past cases \\
\textbf{Slice thickness buttons} & Switch 3\,mm / TS / other; COR/SAG reconstruct from axial if needed \\
\textbf{Finding type} & PM nodule (R) default; also ascites, omental cake, stranding, lymph node \\
\bottomrule
\end{tabular}
\end{center}

\vfill
\hrule
\vspace{0.12cm}
{\footnotesize\color{gray} Outputs: \texttt{<Output>/<case\_id>/R\_1.mrk.json} (or A/O/S/L) with summary in \texttt{\_annotation\_done.json}.}

\end{document}
